\chapter{$\Wass_1$ Optimal Transport}
\label{c-w1}

本章では、ground cost が距離に等しい optimal transport 問題に焦点を当てる。歴史的に、これは\citeauthor{Monge1781}が\citeyear{Monge1781}年に提起した元々の問題に対応する。この設定は、コンピュータビジョンにおける optimal transport の初期の応用~\citep{RubTomGui00}においても「earth mover's distances」の名前で採用されたものである。

ground cost が\emph{二乗} Hilbert 距離である場合（特に第\ref{c-dynamic}章で研究される）とは異なり、コストが距離である輸送問題は理論的に解析がより困難である。
注意~\ref{rem-exist-mongemap}では、二乗距離を用いた場合に2つの絶対連続な測度間の輸送写像またはカップリングの一意性が述べられているが、コストが ground 距離そのものである場合、最適な Kantorovich カップリングは一般に一意ではない。したがって、この体制では一意に定まる Monge 写像を復元することがしばしば不可能であり、この種の問題は測度の補間には不向きである。
最適な $\Wass_1$ 輸送計画の存在証明とその幾何学的構造の詳細な解析については、\citet{trudinger2001monge,caffarelli2002constructing,sudakov1979geometric,evans1999differential}の研究を参照されたい。

理論的な解析はより困難であるものの、線形の ground 距離を用いた optimal transport は、二次コストと比較して外れ値やノイズに対して通常より頑健である。さらに、コストが距離であることにより、局所的なフロー、発散制約、または双対ポテンシャルの Lipschitz 性を含む洗練された双対再定式化が可能となり、グラフ理論における\emph{最小費用フロー}法に適合するより安価な数値アルゴリズムが示唆される。
この設定はまた、関連する OT 距離が任意の分布（正でなくてもよい）を比較できるノルムを定義するため人気がある。この性質は、いわゆる\emph{双対ノルム}のより広いクラスにも共有される（詳細は第\ref{sec-dual-norms}節および注意~\ref{rem-unb-dualnorms}を参照）。

%%%%%%%%%%%%%%%%%%
\section{距離空間上の $\Wass_1$}
\label{sec-w1-metric}

ここでは $d$ が $\Xx=\Yy$ 上の距離であると仮定し、ground cost $c(x,y)=d(x,y)$ を用いて OT 問題を解く。以下の命題は、この設定における $c$-変換~\eqref{eq-c-transform}の重要な性質を明らかにする。以下では、関数 $f \in \Cc(\Xx)$ の Lipschitz 定数を次のように表す：
\eq{
	\Lip(f) \eqdef \sup \enscond{ \frac{|\f(x)-\f(y)|}{d(x,y)} }{ (x,y) \in \X^2, x \neq y }.
}
$\Lip(f) < +\infty$ を満たす関数 $f$ を Lipschitz 関数と定義する。これらは $\Cc(\X)$ の凸部分集合をなす。

\begin{prop}$\X=\Y$ かつ $c(x,y)=d(x,y)$ と仮定する。
このとき、ある $g$ が存在して $\f = g^c$ が成り立つための必要十分条件は $\Lip(f) \leq 1$ である。
さらに、$\Lip(f)\leq1$ ならば $f^c=-f$ である。
\end{prop}

\begin{proof}
まず、$f=g^c$ と仮定する。$x,y\in \X$ に対して、
\begin{align*}
|f(x)-f(y)|&=\left|\inf_{z\in\X} d(x,z)-g(z) \;\; - \;\;\inf_{z\in \X} d(y,z)-g(z) \right|\\
&\leq \sup_{z\in\X} |d(x,z)-d(y,z)|\leq d(x,y).
\end{align*}
最初の等号は $g^c$ の定義から、次の不等式は恒等式 $|\inf f-\inf g|\leq\sup|f-g|$ から、最後の不等式は三角不等式から従う。これにより $\Lip(f)\leq1$ が示される。

次に、$\Lip(f)\leq1$ と仮定し、$g\eqdef -f$ と定義する。Lipschitz 性により、すべての $x,y\in\X$ に対して
$f(y)-d(x,y)\leq f(x)\leq f(y)+d(x,y)$ が成り立つ。
これらの不等式を適用すると、
\begin{align*}
g^c(y)&=\inf_{x\in\X}\left[ d(x,y) + f(x) \right]
\geq\inf_{x\in \X} \left[  d(x,y)+f(y)-d(x,y)\right]=f(y),\\
g^c(y)&=\inf_{x\in\X}\left[ d(x,y) + f(x) \right]
\leq \inf_{x\in\X}\left[ d(x,y) + f(y)+d(x,y) \right]=f(y).
\end{align*}
ゆえに、$g=-f$ として $f=g^c$ が成り立つ。同じ不等式を用いると、
\begin{align*}
f^c(y)&=\inf_{x\in\X}\left[ d(x,y) - f(x) \right]
\geq\inf_{x\in \X} \left[  d(x,y)-f(y)-d(x,y)\right]=-f(y),\\
f^c(y)&=\inf_{x\in\X}\left[ d(x,y) - f(x) \right]
\leq \inf_{x\in\X}\left[ d(x,y) - f(y)+d(x,y) \right]=-f(y).
\end{align*}
これにより $f^c=-f$ が示される。
\end{proof}

単一ポテンシャル定式化~\eqref{eq-semi-dual-cont}から出発し、構成を反復して組 $(\g,\g^c)$ を $(\g^c,(\g^c)^c)$ に置き換えることができる。前の命題により $(\g^c,-\g^c)$ を用いることができ、これは $\Lip(f) \leq 1$ を満たす任意の組 $(f,-f)$ と等価である。これにより、$\Wass_1$ 距離の次の代替的な表現が得られる：
\eql{\label{eq-w1-metric}
	\Wass_1(\al,\be) =
	\umax{f} \enscond{ \int_{\Xx} f(x) (\d\al(x)-\d\be(x)) }{ \Lip(f) \leq 1 }.
}
この表現は、$\Wass_1$ が実際にはノルムであること、\ie $\Wass_1(\al,\be) = \norm{\al-\be}_{\Wass_1}$ であることを示しており、$\int_{\Xx} \al=\int_{\Xx} \be$ を満たす限り、任意の測度（必ずしも正でなくてよい）に対しても有効である。このノルムはしばしば\citeauthor{kantorovich1958space}ノルム~\citeyearpar{kantorovich1958space}と呼ばれる。

式~\eqref{eq-discr-meas}の形の離散測度に対し、$\al-\be = \sum_{k} \VectMode{m}_k \de_{z_k}$（$z_k \in \Xx$、$\sum_k \VectMode{m}_k=0$）と書くと、最適化~\eqref{eq-w1-metric}は次のように書き換えられる：
\eql{\label{eq-w1-discr}
	\Wass_1(\al,\be) =
	\umax{ (\fD_k)_k } \enscond{ \sum_k \fD_k \VectMode{m}_k }{ \foralls (k,\ell),
		\abs{\fD_k-\fD_\ell} \leq d(z_k,z_\ell), }
}
これは二次錐制約を持つ有限次元凸計画問題である。内点法を用いて解くことができるほか、次に類似の問題について詳述するように近接法を用いて解くこともできる。

$\X=\RR$ 上で $d(x,y)=|x-y|$ を用いる場合、$z_k$ を $z_1 \leq z_2 \leq \ldots$ と順序付けることで制約の数を減らすことができる。この場合、解くべき問題は
\eq{
	\Wass_1(\al,\be) =
	\umax{ (\fD_k)_k } \enscond{ \sum_k \fD_k \VectMode{m}_k }{ \foralls k, \abs{\fD_{k+1}-\fD_k} \leq z_{k+1}-z_k },
}
となり、これは線形計画問題である。
なお、この1次元の場合には、累積関数を用いた $\Wass_1$ の閉形式表現が式~\eqref{eq-w1-1d}で与えられる。

\begin{rem}[$0<p \leq 1$ の場合の $\Wass_p$]
	$0<p \leq 1$ ならば、$\tilde d(x,y) \eqdef d(x,y)^p$ は三角不等式を満たし、したがって $\tilde d$ 自体が距離となる。よって、$d$ の代わりに $\tilde d$ を用いることで、上述の $\Wass_1$ に関する結果とアルゴリズムを $\Wass_p$ の計算に適用できる。これは、$\Wass_p$ が $p$-H\"older 関数 $\ensconds{f}{\Lip_p(f) \leq 1}$ の双対であるという主張と等価であり、ここで
	\eq{
		\Lip_p(f) \eqdef \sup \enscond{ \frac{|\f(x)-\f(y)|}{d(x,y)^p} }{ (x,y) \in \X^2, x \neq y }
	}
	である。
\end{rem}

%%%%%%%%%%%%%%%%%%
\section{ユークリッド空間上の $\Wass_1$}
\label{sec-w1-eucl}


ユークリッド空間 $\X=\Y=\RR^\dim$ の特殊な場合に、$\c(x,y) = \norm{x-y}$ を用いると、式~\eqref{eq-w1-metric}に現れる大域的な Lipschitz 制約は $f$ の勾配に対する一様な上界として局所化でき、
\eql{\label{eq-w1-cont}
	\Wass_1(\al,\be) =
	\umax{f} \enscond{ \int_{\RR^\dim} \f(x) (\d\al(x)-\d\be(x)) }{ \norm{\nabla \f}_\infty \leq 1 }
}
となる。ここで制約 $\norm{\nabla \f}_\infty \leq 1$ は、$\f$ の勾配のノルムが任意の点 $x$ において $1$ 以下であること、すなわちすべての $x$ に対して $\norm{\nabla \f(x)}_2 \leq 1$ であることを意味する。

式~\eqref{eq-w1-cont}の双対問題を考えると、固定された発散制約の下での最適化問題が得られる：
\eql{\label{eq-w1-cont-div}
	\Wass_1(\al,\be) =
	\umin{\flow} \enscond{ \int_{\RR^\dim} \norm{\flow(x)}_2 \d x }{  \diverg(\flow)=\al-\be },
}
これはしばしば Beckmann 定式化~\citep{Beckmann52}と呼ばれる。
ここでベクトル値関数 $\flow(x) \in \RR^2$ はフロー場と解釈でき、質量の局所的な移動を記述する。2つの入力測度の台の外側では $\diverg(\flow)=0$ であり、これは質量保存の制約である。
有限要素を用いて適切に離散化すると、問題~\eqref{eq-w1-cont}および\eqref{eq-w1-cont-div}は非平滑凸最適化問題となる。
既成の内点法二次錐最適化ソルバーを使用することも可能であるが、第\ref{sec-prox-solvers}節で提唱されるように、大規模問題にはより単純でありながらより適応的な一次法の使用が必要である。例えば、Douglas--Rachford (DR) 反復~\eqref{eq-dr-iters}や関連する交互方向乗数法を用いることができる。一様グリッド上では、発散制約への射影は高速 Fourier 変換を用いて便利に処理できることに注意されたい。これらのアプローチの詳細な説明および三角形メッシュ上の OT への応用については\citet{SolomonEMDSurfaces2014}を参照されたい。
主双対分離スキームを用いた類似のアプローチについては\citet{LiRyuOsherYinGangbo2017_parallel,Ryu2017a,Ryu2017b}も参照されたい。
双対ポテンシャル $\f$ に対する Lipschitz 制約を緩和する近似スキームも提案されている。例えば、ウェーブレット係数に対する制約を用いて陽的な公式を導く手法~\citep{shirdhonkar2008approximate}や、$\f$ を「rectified linear」$\max(0,\cdot)$ 活性化関数とクリップされた重みを持つ多層ニューラルネットワークとしてパラメトライズする手法~\citep{WassersteinGAN}がある。

%%%%%%%%%%%%%%%%%%
\section{グラフ上の $\Wass_1$}

$\Wass_1$ の前述の定式化~\eqref{eq-w1-cont}および\eqref{eq-w1-cont-div}は、$\X$ が測地線空間である設定、\ie $\c(x,y)=\dist(x,y)$（$\dist$ は測地線距離）に一般化できる。$\X$ が Riemann 多様体である場合の理論的解析については\citet{FeldmanMacCann02}を参照されたい。
$\X=\range{1,n}$ が離散集合であり、重み（長さ）$\weight_{i,j}$ でラベル付けされた無向辺 $(i,j) \in \Edges \subset \X^2$ が備わっている場合、$\X$ が測地線距離（最短経路距離）を備えたグラフであるという重要な場合に帰着する：
\eq{
	\distD_{i,j} \eqdef \umin{K \geq 0, (i_k)_k : i \rightarrow j }
	 \enscond{ \sum_{k=1}^{K-1} \weight_{i_k,i_{k+1}} }{ \foralls k \in \range{1,K-1}, (i_k,i_{k+1}) \in \Ee  },
}
ここで $i \rightarrow j$ は $i_1=i$ かつ $i_K=j$ であること、すなわち経路が $i$ で始まり $j$ で終わることを意味する。

グラフ $\X$ 上の（符号付き）離散測度を定義する2つのベクトル $(\a,\b) \in (\RR^n)^2$ であって $\sum_i \a_i = \sum_i \b_i$ を満たすもの（これらの重みは正である必要はない）を考える。
目標は、ground 距離がグラフの測地線距離である場合に、$p=1$ として式~\eqref{eq-wass-p-disc}で導入された $\WassD_1(\a,\b)$ を計算することである。この計算は、$n \times n$ サイズの「完全な」カップリング $\P$ を計算するところまで行く必要なく、代わりにグラフの基礎的な連結性を活用した局所的な演算子に依拠すべきである。これらの演算子は、勾配および発散微分演算子の離散的定式化である。

離散双対 Kantorovich ポテンシャル $\fD \in \RR^n$ はグラフのすべての頂点で添字付けられたベクトルである。勾配演算子 $\nabla : \RR^n \rightarrow \RR^{\Edges}$ は次のように定義される：
\eq{
	\foralls (i,j) \in \Edges, \quad (\nabla \fD)_{i,j} \eqdef \fD_i-\fD_j.
}
フロー $\flowD = (\flowD_{i,j})_{i,j}$ は辺上で定義され、発散演算子 $\diverg : \RR^{\Edges} \rightarrow \RR^n$ は勾配 $\nabla$ の随伴であり、フローを頂点上で定義されたベクトルに写像し、次のように定義される：
\eq{
	\foralls i \in \range{1,n}, \quad \diverg(\flowD)_i \eqdef \sum_{j: (i,j) \in \Edges} (\flowD_{i,j} - \flowD_{j,i}) \in \RR^n.
}


問題~\eqref{eq-w1-cont}はグラフの設定では次のようになる：
\eql{\label{eq-w1-graph-potential}
	\WassD_1(\a,\b) =
	\umax{\fD \in \RR^n} \enscond{ \sum_{i=1}^n \fD_i (\a_i-\b_i)   }{  \forall (i,j) \in \Edges, \, |(\nabla \fD)_{i,j}| \leq \weight_{i,j}  }.
}
関連する双対問題は、式~\eqref{eq-w1-cont-div}に類似して、次のようになる：
\eql{\label{eq-w1-graph-div}
	\WassD_1(\a,\b) =
	\umin{\flowD \in \RR_+^\Edges} \enscond{ \sum_{(i,j) \in \Ee} \weight_{i,j} \flowD_{i,j}  }{  \diverg(\flowD)=\a-\b }.
}
これは線形計画問題であり、より正確には最小費用フロー問題の一例である。これを解くための非常に効率的な専用シンプレックスソルバーが考案されている。例えば\citep{TreeEMD2007}を参照されたい。
図~\ref{fig-planar-graph}に主問題と双対問題の解の例を示す。
定式化~\eqref{eq-w1-graph-div}はいわゆる Beckmann 定式化~\citep{Beckmann52}であり、交通渋滞モデルの定義と研究に使用・拡張されてきた。例えば\citep{CarlierSantambrogioOTCongestion2008}を参照されたい。



\begin{figure}[h!]
\centering
\begin{tabular}{@{}c@{\hspace{10mm}}c@{}}
\includegraphics[width=.45\linewidth]{w1-graph/potential}&
\includegraphics[width=.45\linewidth]{w1-graph/flow}\\
$\fD$ & $({\color{red}\a},{\color{blue}\b})$ and $\flowD$
\end{tabular}
\caption{\label{fig-planar-graph}
一様な重み $\weight_{i,j}=1$ を持つ平面グラフ上での $\WassD_1(\a,\b)$ の計算例。
左：式~\eqref{eq-w1-graph-potential}の解であるポテンシャル $\fD$（赤から青に向かって値が増加）。辺の緑色の濃さは $|(\nabla \fD)_{i,j}|$ に比例する。
右：式~\eqref{eq-w1-graph-div}の解であるフロー $\flowD$。太い黒の辺は非零の $\flowD_{i,j}$ を表し、$\weight_{i,j}=1$ に飽和している。
右の飽和したフロー辺は、左の $|(\nabla \fD)_{i,j}|=1$ である明るい緑の辺に対応する。
}
\end{figure}
